% GENERATED FILE. Edit canonical lessons, then run npm run generate:books.
% canonical-source-hash: fnv1a64:2bb43bcc5d911a5f
% canonical-lessons: SA-C279-bhavan-gacchet, SA-C279-atah-tasmat, SA-C279-upadesah, SA-R279-advice-and-what-follows

\chapter{Advice, and What Follows}
\label{ch:sa-advice-and-what-follows}
\hlchaptermodality{279}

\begin{chapteropening}
\textbf{By the end of this chapter:} \emph{I can give advice (bhavān vaidyaṃ gacchet, tvaṃ paṭheḥ) and say what follows from a reason (ataḥ, tasmāt).}

Complete SA-R279-advice-and-what-follows: give advice and say what follows from a reason.
\end{chapteropening}

\section[\texorpdfstring{\sk{भवान्} \sk{गच्छेत्} (bhavān gacchet)}{भवान् गच्छेत् (bhavān gacchet)}]{\sk{भवान्} \sk{गच्छेत्} (bhavān gacchet) — you should go}

\label{lesson:SA-C279-bhavan-gacchet}



\begin{quote}
Say \textquotedblleft{}I will not come, because I am ill\textquotedblright{}. (\emph{ahaṃ na āgamiṣyāmi, yataḥ ahaṃ rugṇaḥ asmi}.) A friend answers with advice.
\end{quote}

\subsection*{You'll want to know: -\sk{एत्}, \textquotedblleft{}should\textquotedblright{}}
\begin{quote}
\textbf{\sk{भवान्} \sk{गच्छेत्।}}
\end{quote}

\begin{quote}
\emph{bhavān gacchet.} — \textquotedblleft{}You should go.\textquotedblright{}
\end{quote}

\textbf{\sk{गच्छेत्}} (\emph{gacchet}) is \textbf{\sk{गच्छति}} with \emph{-et} in place of \emph{-ati}: \textquotedblleft{}should go\textquotedblright{}. It speaks to \textbf{\sk{भवान्}} (\emph{bhavān}), \textquotedblleft{}you\textquotedblright{}, said with respect, and about anyone else:

\begin{quote}
\textbf{\sk{भवान्} \sk{वैद्यं} \sk{गच्छेत्।}}
\end{quote}

\begin{quote}
\emph{bhavān vaidyaṃ gacchet.} — \textquotedblleft{}You should go to the doctor.\textquotedblright{}
\end{quote}

\begin{quote}
\textbf{\sk{सः} \sk{सत्यं} \sk{वदेत्।}}
\end{quote}

\begin{quote}
\emph{saḥ satyaṃ vadet.} — \textquotedblleft{}He should tell the truth.\textquotedblright{}
\end{quote}

To a woman, \textbf{\sk{भवती}} (\emph{bhavatī}):

\begin{quote}
\textbf{\sk{भवती} \sk{गृहे} \sk{तिष्ठेत्।}}
\end{quote}

\begin{quote}
\emph{bhavatī gṛhe tiṣṭhet.} — \textquotedblleft{}You should stay at home.\textquotedblright{}
\end{quote}

\textbf{\sk{भवेत्}} (\emph{bhavet}), from \textbf{\sk{भवति}}, is \textquotedblleft{}should be\textquotedblright{}, and also \textquotedblleft{}may be\textquotedblright{}:

\begin{quote}
\textbf{\sk{सः} \sk{गृहे} \sk{भवेत्।}}
\end{quote}

\begin{quote}
\emph{saḥ gṛhe bhavet.} — \textquotedblleft{}He may be at home.\textquotedblright{}
\end{quote}

Pāṇini (3.3.161) gives this form for bidding and asking.

\subsection*{You'll want to know: \sk{त्वं} \sk{पठेः}, and \sk{मया} \sk{गन्तव्यम्}}
To a friend, with \textbf{\sk{त्वम्}}, the ending is \emph{-eḥ}:

\begin{quote}
\textbf{\sk{त्वं} \sk{प्रतिदिनं} \sk{पठेः।}}
\end{quote}

\begin{quote}
\emph{tvaṃ pratidinaṃ paṭheḥ.} — \textquotedblleft{}You should read every day.\textquotedblright{}
\end{quote}

For yourself, \textquotedblleft{}I must\textquotedblright{} can be said another way, word for word \textquotedblleft{}by me, to be gone\textquotedblright{}:

\begin{quote}
\textbf{\sk{मया} \sk{गन्तव्यम्।}}
\end{quote}

\begin{quote}
\emph{mayā gantavyam.} — \textquotedblleft{}I must go.\textquotedblright{}
\end{quote}

\textbf{\sk{गन्तव्यम्}} (\emph{gantavyam}) is shaped like \textbf{\sk{कर्तव्यम्}} (\emph{kartavyam}), \textquotedblleft{}a duty\textquotedblright{}, and \textbf{\sk{मया}} (\emph{mayā}) is \textquotedblleft{}by me\textquotedblright{}. Meet it when you hear it; you do not need to build it.

\noindent
\begin{tabularx}{\linewidth}{@{}>{\raggedright\arraybackslash}X>{\raggedright\arraybackslash}X@{}}
\toprule
\textbf{} & \textbf{} \\
\midrule
\textbf{\sk{भवान्} \sk{गच्छेत्}} & you should go \\
\textbf{\sk{त्वं} \sk{पठेः}} & you (a friend) should read \\
\textbf{\sk{सः} \sk{वदेत्}} & he should say \\
\textbf{\sk{मया} \sk{गन्तव्यम्}} & I must go \\
\bottomrule
\end{tabularx}

These verbs end in \emph{-ati}, like \textbf{\sk{गच्छति}}; a verb such as \textbf{\sk{करोति}} has an optative of its own, met one at a time.

\subsection*{Your turn}
Say these aloud:

\begin{itemize}
\raggedright
  \item \textquotedblleft{}You should go home\textquotedblright{} — \emph{bhavān gṛhaṃ gacchet}
  \item \textquotedblleft{}You should stay at home\textquotedblright{} — \emph{bhavān gṛhe tiṣṭhet}
  \item \textquotedblleft{}You should read at home\textquotedblright{} — \emph{bhavān gṛhe paṭhet}
\end{itemize}

\begin{tcolorbox}[breakable,colback=teal!4,colframe=teal!35!black,title={Before you move on}]
\textquotedblleft{}You should go\textquotedblright{}, with respect? (\textbf{\sk{भवान्} \sk{गच्छेत्}}.) \textquotedblleft{}You should read\textquotedblright{}, to a friend? (\textbf{\sk{त्वं} \sk{पठेः}}.) \textquotedblleft{}He may be at home\textquotedblright{}? (\textbf{\sk{सः} \sk{गृहे} \sk{भवेत्}}.) What does \textbf{\sk{मया} \sk{गन्तव्यम्}} say? (\textquotedblleft{}I must go\textquotedblright{}.)
\end{tcolorbox}
\section[\texorpdfstring{\sk{अतः} · \sk{तस्मात्} (ataḥ · tasmāt)}{अतः · तस्मात् (ataḥ · tasmāt)}]{\sk{अतः} · \sk{तस्मात्} (ataḥ · tasmāt) — so, and therefore}

\label{lesson:SA-C279-atah-tasmat}



\begin{quote}
Say \textquotedblleft{}you should go to the doctor\textquotedblright{}. (\emph{bhavān vaidyaṃ gacchet}.) Now say why, and what follows.
\end{quote}

\subsection*{You'll want to know: \sk{अतः}}
\textbf{\sk{यतः}} gives the reason after the event; \textbf{\sk{अतः}} (\emph{ataḥ}), \textquotedblleft{}so\textquotedblright{}, gives the result after the reason:

\begin{quote}
\textbf{\sk{अहं} \sk{रुग्णः} \sk{अस्मि।}}
\end{quote}

\begin{quote}
\emph{ahaṃ rugṇaḥ asmi.} — \textquotedblleft{}I am ill.\textquotedblright{}
\end{quote}

\begin{quote}
\textbf{\sk{अतः} \sk{अहं} \sk{न} \sk{आगमिष्यामि।}}
\end{quote}

\begin{quote}
\emph{ataḥ ahaṃ na āgamiṣyāmi.} — \textquotedblleft{}So I will not come.\textquotedblright{}
\end{quote}

\begin{quote}
\textbf{\sk{वर्षा} \sk{भवति}, \sk{अतः} \sk{सः} \sk{गृहे} \sk{अस्ति।}}
\end{quote}

\begin{quote}
\emph{varṣā bhavati, ataḥ saḥ gṛhe asti.} — \textquotedblleft{}It is raining, so he is at home.\textquotedblright{}
\end{quote}

\subsection*{You'll want to know: \sk{तस्मात्}, \sk{ततः}}
\textbf{\sk{तस्मात्}} (\emph{tasmāt}), \textquotedblleft{}from that\textquotedblright{}, says the same with more weight: \textquotedblleft{}therefore\textquotedblright{}. \textbf{\sk{ततः}} (\emph{tataḥ}), \textquotedblleft{}then\textquotedblright{} in a story, can also mean \textquotedblleft{}so\textquotedblright{}:

\begin{quote}
\textbf{\sk{सा} \sk{श्रान्ता} \sk{अस्ति।}}
\end{quote}

\begin{quote}
\emph{sā śrāntā asti.} — \textquotedblleft{}She is tired.\textquotedblright{}
\end{quote}

\begin{quote}
\textbf{\sk{तस्मात्} \sk{सा} \sk{गृहे} \sk{तिष्ठति।}}
\end{quote}

\begin{quote}
\emph{tasmāt sā gṛhe tiṣṭhati.} — \textquotedblleft{}Therefore she is staying at home.\textquotedblright{}
\end{quote}

\begin{quote}
\textbf{\sk{धूमयानं} \sk{न} \sk{आगतम्}, \sk{ततः} \sk{वयं} \sk{गतवन्तः।}}
\end{quote}

\begin{quote}
\emph{dhūmayānaṃ na āgatam, tataḥ vayaṃ gatavantaḥ.} — \textquotedblleft{}The train did not come, so we left.\textquotedblright{}
\end{quote}

\textbf{\sk{आगतम्}} (\emph{āgatam}) is \textbf{\sk{आगतः}}, \textquotedblleft{}come\textquotedblright{}, for a neuter noun such as \textbf{\sk{धूमयानम्}}.

\subsection*{Your turn}
Say these aloud:

\begin{itemize}
\raggedright
  \item \textquotedblleft{}So I will stay at home\textquotedblright{} — \emph{ataḥ ahaṃ gṛhe sthāsyāmi}
  \item \textquotedblleft{}So I will read at home\textquotedblright{} — \emph{ataḥ ahaṃ gṛhe paṭhiṣyāmi}
  \item \textquotedblleft{}So I will eat at home\textquotedblright{} — \emph{ataḥ ahaṃ gṛhe khādiṣyāmi}
\end{itemize}

\begin{tcolorbox}[breakable,colback=teal!4,colframe=teal!35!black,title={Before you move on}]
Which word opens the reason, \textbf{\sk{यतः}} or \textbf{\sk{अतः}}? (\textbf{\sk{यतः}}.) And the result? (\textbf{\sk{अतः}}, or \textbf{\sk{तस्मात्}}.) \textquotedblleft{}She is tired\textquotedblright{}? (\textbf{\sk{सा} \sk{श्रान्ता} \sk{अस्ति}}.)
\end{tcolorbox}
\section[\texorpdfstring{\sk{उपदेशः} (upadeśaḥ)}{उपदेशः (upadeśaḥ)}]{\sk{उपदेशः} (upadeśaḥ) — a piece of advice}

\label{lesson:SA-C279-upadesah}



\begin{quote}
Before the new one: say \textquotedblleft{}you should go to the doctor\textquotedblright{}. (\emph{bhavān vaidyaṃ gacchet}.) That sentence is a piece of advice.
\end{quote}

\subsection*{You'll want to know: \sk{उपदेशः}}
\textbf{\sk{उपदेशः}} — \emph{upadeśaḥ} — \textquotedblleft{}a piece of advice\textquotedblright{}. It is a masculine noun.

\begin{quote}
\textbf{\sk{मम} \sk{उपदेशः} \sk{अस्ति।}}
\end{quote}

\begin{quote}
\emph{mama upadeśaḥ asti.} — \textquotedblleft{}I have a piece of advice.\textquotedblright{}
\end{quote}

\begin{quote}
\textbf{\sk{सः} \sk{मम} \sk{उपदेशं} \sk{न} \sk{शृणोति।}}
\end{quote}

\begin{quote}
\emph{saḥ mama upadeśaṃ na śṛṇoti.} — \textquotedblleft{}He does not listen to my advice.\textquotedblright{}
\end{quote}

\begin{quote}
\textbf{\sk{वैद्यः} \sk{उपदेशं} \sk{दास्यति।}}
\end{quote}

\begin{quote}
\emph{vaidyaḥ upadeśaṃ dāsyati.} — \textquotedblleft{}The doctor will give advice.\textquotedblright{}
\end{quote}

\begin{cousinweb}
\textbf{\sk{उपदेशः}} (\emph{upadeśaḥ}) is built from \textbf{\sk{उप}} (\emph{upa}), \textquotedblleft{}towards\textquotedblright{}, and the root √diś, \textquotedblleft{}to point out\textquotedblright{}. Monier-Williams gives its first sense as pointing something out to someone, and then instruction, teaching and advice.
\end{cousinweb}

\subsection*{Your turn}
Say these aloud:

\begin{itemize}
\raggedright
  \item \emph{upadeśaḥ}
  \item \textquotedblleft{}I have a piece of advice\textquotedblright{} — \emph{mama upadeśaḥ asti}
  \item \textquotedblleft{}My advice: you should rest at home\textquotedblright{} — \emph{mama upadeśaḥ: bhavān gṛhe tiṣṭhet}
\end{itemize}

\begin{tcolorbox}[breakable,colback=teal!4,colframe=teal!35!black,title={Before you move on}]
What does \textbf{\sk{उपदेशः}} mean? (\textquotedblleft{}A piece of advice\textquotedblright{}.) Masculine or neuter? (Masculine.)
\end{tcolorbox}
\section[(review)]{Review — advice, and what follows}

\label{lesson:SA-R279-advice-and-what-follows}



\begin{quote}
Say \textquotedblleft{}in my opinion this is beautiful\textquotedblright{}. (\emph{mama matena etat sundaram}.) Then \textquotedblleft{}because of the rain\textquotedblright{}. (\emph{varṣākāraṇāt}.)
\end{quote}

\subsection*{Your turn}
Say these aloud:

\begin{itemize}
\raggedright
  \item You are ill. You should go to the doctor.
  \item You should read every day, to a friend.
  \item It is raining, so we will stay at home.
  \item She is tired, therefore she is not coming.
  \item I have a piece of advice: he should tell the truth.
\end{itemize}

\begin{tcolorbox}[breakable,colback=teal!4,colframe=teal!35!black,title={Before you move on}]
\textbf{\sk{गच्छेत्}} or \textbf{\sk{गच्छति}}, for advice? (\textbf{\sk{गच्छेत्}}.) \textbf{\sk{पठेः}}: to whom? (A friend, \textbf{\sk{त्वम्}}.) Now the other way round: \textquotedblleft{}so, therefore\textquotedblright{}? (\textbf{\sk{अतः}}, \textbf{\sk{तस्मात्}}.) \textquotedblleft{}I must go\textquotedblright{}? (\textbf{\sk{मया} \sk{गन्तव्यम्}}.)
\end{tcolorbox}
